\documentclass[10pt,a4paper]{amsart}
\usepackage[margin=19mm]{geometry}
\usepackage{amsmath,amssymb,mathtools}
\usepackage[scr=rsfs,cal=euler]{mathalfa}
\usepackage{xfrac}
\usepackage[unicode,hidelinks,bookmarks=false]{hyperref}
\newcommand{\ie}{i.e.\ }
\newcommand{\cf}{cf.\ }
\newcommand{\resp}{respectively\ }
\newcommand{\Subsection}{\subsection}
\providecommand{\nobreakdash}{\nobreak\hbox{-}}
\setlength{\emergencystretch}{2em}
\multlinegap0pt
\usepackage{bm}
\usepackage{bbm}
\usepackage{dsfont}
\usepackage{xspace}
\usepackage{cancel}
\usepackage{mathtools}
\usepackage{amsthm}
\makeatletter
\@ifundefined{theorem}{\newtheorem{theorem}{Theorem}}{}
\makeatother
\newcommand{\bR}{{\mathbb R}}
\newcommand{\cO}{{\mathcal O}}
\newcommand{\s}{\sigma}
\title[Quantum Hamming Metrics]{Quantum Hamming Metrics}
\author{Marc A. Rieffel}
\date{Source: arXiv:2507.23046v1 (CC BY 4.0)}
\begin{document}
\maketitle

\noindent\emph{Source and licence.} Quantum Hamming Metrics. Version arXiv:2507.23046v1 (CC BY 4.0), distributed under the Creative Commons Attribution 4.0 International licence, as recorded in the arXiv licence metadata for this exact version and shown on the abstract page https://arxiv.org/abs/2507.23046v1. Full rights notice: licenses/arxiv-2507.23046-CC-BY-4.0.txt.\par\smallskip

\noindent\textit{Editorial scope.} Counted unit: the Theorem whose source label is \texttt{predual} in the article (source0.tex lines 841--859, source label \texttt{predual}), delivered together with the article's own complete original proof of it (source0.tex lines 861--938, one \texttt{proof} environment of 78 lines). No mathematics is written, repaired or re-derived: statement and proof are the article's own text line for line. Required context retained uncounted: none. Every definition, display and label of those lines is retained unchanged. Omitted from this package and not redistributed: 1--840, 860--860, 939--2341. Nothing inside a retained excerpt is omitted or altered: every retained body is delivered exactly as the article's lines are written, so no omission receipt of any kind accompanies this package. Citations: the retained excerpts carry no citation command at all, so no bibliography entry of the article is transcribed and no reference is redistributed. Reference adaptation: none was needed and none was made; every reference inside the delivered unit resolves inside it, so no unresolved reference or citation is produced. Printed slips kept literal: no printed word, symbol, hypothesis or number of the counted statement or of its proof was altered. Layout and class: the article's own class is \texttt{amsart} with packages including amssymb, amsfonts, amsmath, hyperref, cite; the wrapper is the lane's pinned \texttt{amsart} editorial wrapper at 10pt on A4, which restates the theorem-like environments the retained text needs and adds the article's own macro definitions that the retained text needs (\texttt{\textbackslash bR}, \texttt{\textbackslash cO}, \texttt{\textbackslash s}), with extra wrapper packages (bm, bbm, dsfont, xspace, cancel, mathtools). The delivered numbering is the unit's own. Assets: no retained span contains a figure environment, a graphics-inclusion command, a PDF-page inclusion, a table or a TikZ picture; no image file is copied into this package, the package declares no asset dependency and no image file is redistributed with it. Licence: version arXiv:2507.23046v1 of this article is distributed under the Creative Commons Attribution 4.0 International licence, as recorded in the arXiv licence metadata for this exact version and shown on the abstract page https://arxiv.org/abs/2507.23046v1. A full rights notice is delivered at licenses/arxiv-2507.23046-CC-BY-4.0.txt. Characters: every retained excerpt is pure ASCII, so the delivered engine, XeLaTeX with its default Latin Modern text font, reports no missing character.
\begin{theorem}
\label{predual}
Let $\{V, \|\cdot\|_V\}$ be a normed vector space, and let $W$ be a dense 
subspace of $V$. Let $\| \cdot \|_W$ be a norm on $W $ for which there is a 
constant, $r \in \bR$, such that $\|w\|_V \leq r \|w\|_W$ for all $w \in W$.
Then each $\phi \in V'$ can be viewed as a continuous linear functional on $W$,
and $\|\phi\|_{W'} \leq r \|\phi\|_{V'}$, so that $B_{V'} \subseteq r B_{W'}$. 

Suppose that $B_W$ is totally bounded for $\|\cdot\|_V$. Then
the topology on $B_{V'}$ from the norm $\|\cdot\|_{W'}$ coincides
with the restriction to it of the weak-$*$ topology of $V'$ as the
dual of $V$ (for which $B_{V'}$ is compact). 

Furthermore, let $F$ be the subspace of $W'$ consisting
of those elements whose restriction to $B_W$ is continuous for
its topology from $\|\cdot \|_V$. Then $F$ is a predual for 
the completion of $W$, 
and $V'$ is  dense in $F$ for the norm $\|\cdot\|_{W'}$.
 \end{theorem}

\begin{proof}
We show first that for the situation in the first paragraph, the
 topology on $B_{V'}$ from the norm $\|\cdot\|_{W'}$ is stronger than
 the restriction to it of the weak-$*$ topology of $V'$ as the
dual of $V$. For a given $\phi \in  B_{V'}$ let 
\[
\cO(\phi, \{v_j\}, \epsilon) =\{\psi \in  B_{V'}: |(\phi - \psi)(v_j)| < \epsilon, \ \ 
 j = 1, \dots , n\}  ,
\]
 an arbitrary element of a neighborhood base for $\phi$ for the weak-$*$ topology of $V'$
 restricted to $B_{V'}$.
 We seek a $d \in \bR$ such that the open 
 ball $B_{V'}(\phi, d) = \{\psi \in  B_{V'}: \|\phi - \psi\|_{W'} < d\}$ about $\phi$ is contained in
$\cO(\phi, \{v_j\}, \epsilon) $. Let $M = \max\{\|v_j\|: j = 1,\dots , n\}$, and let
$d = \epsilon/rM$. Then we see that for any $\psi \in B_{V'}(\phi, d)$ and each $j$ we 
have $|(\phi - \psi)(v_j)| < \epsilon\}$, as needed.   

Now suppose that $B_W$ is totally bounded for $\|\cdot\|_V$. We show then that 
 the topology on $B_{V'}$ from the norm $\|\cdot\|_{W'}$ is weaker than
 the restriction to it of the weak-$*$ topology of $V'$ as the
dual of $V$ (so these topologies coincide in view of the previous paragraph).
For a given $\phi \in  B_{V'}$ and a $d > 0$, let 
$B_{W'}(\phi, d) = \{\psi \in  B_{V'}: \|\phi - \psi\|_{W'} < d\}$,  
 an arbitrary element of a neighborhood base for $\phi$ for the norm of $W'$.
 We want to show that this ball contains an open neighborhood of $\phi$
 for the weak-$*$ topology of $V'$. Since $B_W$ is totally bounded for $\|\cdot\|_V$,
 we can find $w_1, \dots, w_n$ in $B_W$ such that the $\|\cdot \|_V$-balls
 about them of
 radius $d/3$ cover $B_W$. Set
 \[
\cO(\phi, \{w_j\}, d/3) =\{\psi \in  B_{V'}: |(\phi - \psi)(w_j)| < d/3, \ \ j = 1, \dots , n\}  ,
\]
a weak-$*$ open neighborhood of $\phi$. Consider $w \in B_W$. There is
a $k$ such that $\|w - w_k\|_V < d/3$. Then for any $\psi \in \cO(\phi, \{w_j\}, d/3)$
we have
\begin{align*}
|(\phi & - \psi)(w)| \\
&\leq |\phi(w) - \phi(w_k)| + |\phi(w_k) - \psi(w_k)| + |\psi(w_k) - \psi(w)|  \\
&\leq \|\phi\|_{V'}\|w - w_k\|_V +  |(\phi - \psi)(w_k)| + \|\psi\|_{V'}\|w_k - w\|_V    < d   .
\end{align*}
Thus $\|\phi - \psi\|_{W'} < d$, so that $\psi \in B_{W'}(\phi, d)$ as needed.

We now show that the completion of $W$ is a dual space, assuming as above
that $B_W$ is totally bounded for $\|\cdot\|_V$. 
When we take the completions of $V$ and $W$, the conditions of the theorem
continue to hold. So we now assume that $V$ and $W$ are complete.
Of course a predual, $F$, for $W$ must
be a closed subspace of $W'$. Let $F$ consist of all the elements of $W'$
whose restriction to $B_W$ is continuous for the (compact) topology on $B_W$ 
from $\|\cdot\|_V$. Then $F$ is a closed subspace of $W'$.
Clearly $V' \subseteq F$, so $F$ separates the points of $W$. Define a linear
map, $\s$,  from $W$ into $F'$ by $\s(w)(f) = f(w)$. Then $\s$ is injective, and
$|\s(w)(f)| \leq \|f\|_{W'}\|w\|_W$, so that $\|\s\| \leq 1$. 
In particular, $\s(B_W) \subseteq B_{F'}$. Since the composition of $\s$
with the linear functional on $F'$ determined by any element of $F$ is, from the 
definition of $F$, obviously
 continuous on $B_W$ for its compact topology from $\|\cdot\|_V$, it follows that $\s$ 
 on $B_W$ is continuous
 for the weak-$*$ topology on $F'$. It follows that $\s(B_W)$ is compact for the
 weak-$*$ topology. 
 
 Since  $\s(B_W)$ is convex, the Hahn-Banach separation
 theorem tells us that
if $\s(B_W)$ were not all of $B_{F'}$ there would be a $\theta \in B_{F'}$ and
a weak-$*$ linear functional separating $\theta$ from $\s(B_W)$.  
But every weak-$*$ linear functional on $F'$ comes from an element of $F$.
So there would be an $f \in F$ such that $f(w) \leq 1 < \theta(f)$ for every
$w \in B_W$. But the first inequality says that $\|f\|_{W'} \leq 1$, and so
the second inequality says that $\|\theta\|_{F'} > 1$, contradicting
$\theta \in B_{F'}$. Thus $\s$ is an isometry from $W$ onto $F'$.

Finally, suppose that there is a $\theta \in F'$ such that 
$\theta(\phi) = 0$ for all $\phi \in V' \subseteq F$.
We have just seen that every element of $F'$ comes from an element of $W$,
so there is a $w \in W$ such that $0= \theta(\phi) =  \phi(w)$ for all $\phi$.
It follows that $w = 0$, so that $\theta = 0$. Then the Hahn-Banach theorem
tells us that the closure of $V'$ in $F'$ for the norm $\|\cdot\|_{W'}$ is all of $F'$.
\end{proof} 

\end{document}
